\providecommand{\cref}[1]{\ref{#1}}
\section{Wall-crossing gluing}
\label{sec:gluing:wallcrossing}
\label{part:gluing-wallcrossing}
\subsection{The Kontsevich--Soibelman quantum torus}
\label{sec:ks-quantum-torus}

\subsubsection{Charge lattice and skew-Euler form}

Fix a CY$_3$ triangulated category $\cC$ with finitely generated
Grothendieck group $\Gamma \defeq K_0(\cC)$. Every object $E \in \cC$
defines a class $[E] \in \Gamma$, and every short exact triangle
$E' \to E \to E''$ gives $[E] = [E'] + [E'']$ in $\Gamma$. The Euler
form is
\[
 \chi(E, F) \;\defeq\; \sum_{i \in \Z} (-1)^i \dim \Ext^i_\cC(E, F),
\]
a bilinear pairing $\Gamma \times \Gamma \to \Z$. Set
\[
 \langle \gamma, \gamma' \rangle \;\defeq\; \chi(\gamma, \gamma') - \chi(\gamma', \gamma).
\]
The antisymmetrisation cancels the symmetric part; what remains is the
\emph{skew-Euler form}. On any triangulated category it is integer-valued and
bilinear in both slots; on a CY$_d$ category Serre duality provides an
identification $\Ext^i(E, F) \cong \Ext^{d-i}(F, E)^\vee$, which forces a
symmetry relation we exploit below.

\begin{definition}[Quantum torus {\cite[\S 2.3]{KS2008}}]
 \label{def:ks-quantum-torus}
 Let $\LL^{1/2}$ be a formal variable (motivic square-root of the Lefschetz
 motive $\LL = \LL^{1/2} \cdot \LL^{1/2}$). The \emph{Kontsevich--Soibelman
 quantum torus} of $(\Gamma, \langle \cdot, \cdot \rangle)$ is the algebra
 \[
 \mathcal{T}_\Gamma \;\defeq\; \Big(\bigoplus_{\gamma \in \Gamma} \Z[\LL^{\pm 1/2}] \,\hat{x}_\gamma \Big),
 \qquad
 \hat{x}_\gamma \cdot \hat{x}_{\gamma'} \;=\; \LL^{\langle \gamma, \gamma' \rangle / 2}\, \hat{x}_{\gamma + \gamma'}.
 \]
 Passing to the appropriate completion over the positive cone
 $\Gamma^+ \subset \Gamma$ gives the \emph{completed quantum torus}
 $\hat{\mathcal{T}}_\Gamma = \prod_{\gamma \in \Gamma^+} \Z[\LL^{\pm 1/2}]\,\hat{x}_\gamma$, which is where
 the wall-crossing formula lives.
\end{definition}

\begin{remark}[Conventions]
 The convention for the exponent varies among sources: Kontsevich--Soibelman
 use $\langle \gamma, \gamma' \rangle / 2$ with $\langle \cdot, \cdot \rangle$
 the skew-Euler form; Joyce--Song use $\chi(\gamma, \gamma')$ (not
 antisymmetrised) with a sign in the Lie bracket; Nagao et al.~use
 $-\chi(\gamma, \gamma')$. All three conventions produce isomorphic quantum
 tori via a gauge change on $\hat{x}_\gamma$. We adopt the KS convention.
\end{remark}

\subsubsection{The skew-Euler form on a CY$_3$ category}

The crucial feature of dimension three is that the Serre duality operator
\emph{makes the Euler form itself antisymmetric up to sign}. Elementary:

\begin{proposition}[CY$_3$ Serre antisymmetry]
 \label{prop:cy3-serre-antisymmetry}
 For $\cC$ a CY$_3$ triangulated category,
 \[
 \chi(\gamma, \gamma') \;=\; -\chi(\gamma', \gamma)
 \qquad \text{for all } \gamma, \gamma' \in \Gamma.
 \]
\end{proposition}

\begin{proof}
 Serre duality on a CY$_d$ category gives
 $\Ext^i(E, F) \cong \Ext^{d-i}(F, E)^\vee$ as $\Z$-modules. Taking alternating
 dimensions,
 \[
 \chi(E, F) = \sum_i (-1)^i \dim \Ext^i(E, F)
 = \sum_i (-1)^i \dim \Ext^{d-i}(F, E)
 = (-1)^d \sum_j (-1)^j \dim \Ext^j(F, E)
 = (-1)^d \chi(F, E).
 \]
 At $d = 3$ the factor is $(-1)^3 = -1$, establishing antisymmetry.
\end{proof}

\begin{corollary}[Skew-Euler equals $2\chi$ on CY$_3$]
 \label{cor:skew-euler-equals-2chi}
 On any CY$_3$ category $\cC$,
 \[
 \langle \gamma, \gamma' \rangle \;=\; \chi(\gamma, \gamma') - \chi(\gamma', \gamma) \;=\;
 \chi(\gamma, \gamma') + \chi(\gamma, \gamma') \;=\; 2\chi(\gamma, \gamma').
 \]
 The Euler form itself is \emph{already} antisymmetric, and the skew-Euler
 form is $2\chi$.
\end{corollary}

Substituting $\langle \gamma, \gamma' \rangle = 2\chi(\gamma, \gamma')$ into
\cref{def:ks-quantum-torus},
\[
 \hat{x}_\gamma \cdot \hat{x}_{\gamma'} \;=\; \LL^{\langle \gamma, \gamma' \rangle / 2}\, \hat{x}_{\gamma + \gamma'}
 \;=\; \LL^{\chi(\gamma, \gamma')}\, \hat{x}_{\gamma + \gamma'}.
\]
The commutator is computed directly from antisymmetry of $\chi$:
\[
 \hat{x}_\gamma \hat{x}_{\gamma'} \;=\; \LL^{\chi(\gamma, \gamma')}\, \hat{x}_{\gamma+\gamma'},
 \qquad
 \hat{x}_{\gamma'} \hat{x}_\gamma \;=\; \LL^{\chi(\gamma', \gamma)}\, \hat{x}_{\gamma+\gamma'}
 \;=\; \LL^{-\chi(\gamma, \gamma')}\, \hat{x}_{\gamma+\gamma'},
\]
so
\[
 \hat{x}_\gamma \hat{x}_{\gamma'} \hat{x}_\gamma^{-1} \hat{x}_{\gamma'}^{-1}
 \;=\; \LL^{\chi(\gamma, \gamma') - \chi(\gamma', \gamma)} \;=\; \LL^{\langle \gamma, \gamma'\rangle}
 \;=\; \LL^{2\chi(\gamma, \gamma')}.
\]
The factor $2$ is a feature of the CY$_3$ skew-Euler form, not a convention.
At the motivic level the quantum torus is non-commutative with this
explicit commutator; in the classical limit $\LL^{1/2} \to 1$, every
power of $\LL$ trivialises and $\hat{x}_\gamma \hat{x}_{\gamma'} = \hat{x}_{\gamma+\gamma'}$ is the
group algebra $\Z[\Gamma]$.

\begin{remark}[$\chi$ is antisymmetric, not symmetric, on CY$_3$]
 \label{rem:chi-antisymmetric-cy3}
 The quantity that is antisymmetric on a CY$_3$ category is $\chi$
 \emph{itself}, by Serre duality through the middle dimension. It is
 \emph{not} the case that one starts with a symmetric form and
 antisymmetrises to obtain the skew-Euler pairing; in lower dimension
 $d$ even, $\chi$ is symmetric, and skew-Euler is a genuine
 antisymmetrisation. At $d = 3$ the symmetrisation
 $\chi(\gamma, \gamma') + \chi(\gamma', \gamma)$ already vanishes, the
 Euler form is already antisymmetric, and the notation $\langle \cdot, \cdot\rangle$
 redundantly duplicates $\chi$ by a factor of two. This is the only
 dimension on which the distinction collapses in this way; the
 quantum-torus commutator $\LL^{2\chi}$ is the motivic shadow of the
 $d=3$ collapse.
\end{remark}

\begin{remark}[Hall-product non-commutativity lives elsewhere]
 \label{rem:hall-product-distinct}
 The non-commutativity $\hat{x}_\gamma \hat{x}_{\gamma'} = \LL^\chi \hat{x}_{\gamma+\gamma'}$ of the quantum
 torus is the \emph{charge-level} non-commutativity: it sees only the
 $K_0$-grading, not the moduli-space geometry. The Hall-product
 non-commutativity on the full critical-cohomology CoHA $\cH(\cC)$ is a
 strictly richer datum, involving the vanishing-cycle complexes on
 moduli stacks of extensions and not merely the Euler pairing of their
 classes. The passage from the quantum torus to the full CoHA is the
 subject of \cref{sec:js-hopf-chamber-independence}; the two
 non-commutativities agree on the $K_0$-graded pieces but disagree on
 finer cohomological data.
\end{remark}

\subsection{Chambers as gluing cells}
\label{sec:chambers-gluing-cells}

\subsubsection{The stability manifold and its stratification}

Bridgeland's stability manifold $\Stab(\cC)$ is a complex manifold carrying
the data needed to specify BPS counts. A point $\sigma \in \Stab(\cC)$
consists of a group homomorphism $Z \colon \Gamma \to \C$ (the
\emph{central charge}) together with a heart $\cA \subset \cC$ of a bounded
$t$-structure, subject to the compatibility that $Z$ maps nonzero objects of
$\cA$ into the upper half plane $\mathbb H \cup \R_{<0}$.

\begin{proposition}[Stab dimension = numerical Grothendieck rank, Bridgeland (2007, arXiv:math/0212237) Thm 7.1]
 \label{prop:stab-dim-equals-K0num-rank}
 For $\cC$ a triangulated category with finite-rank numerical Grothendieck
 group $K_0^{\mathrm{num}}(\cC) \defeq K_0(\cC) / \mathrm{rad}(\chi)$, the
 stability manifold $\Stab(\cC)$ is a complex manifold of complex dimension
 \[
 \dim_\C \Stab(\cC) \;=\; \mathrm{rk}\, K_0^{\mathrm{num}}(\cC).
 \]
 The CY dimension $d$ does \emph{not} enter this formula; the rank is a
 numerical invariant of $\cC$ alone.
\end{proposition}

\begin{example}[Conifold versus strict CY$_3$]
 \label{ex:conifold-strict-cy3-ranks}
 On the resolved conifold $\bar Y$ the numerical Grothendieck group is
 generated by the two simples $\gamma_0, \gamma_1$ of the Klebanov--Witten
 quiver, so $\mathrm{rk}\, K_0^{\mathrm{num}}(\bar Y) = 2$ and $\Stab(\bar
 Y)$ is four-real-dimensional. On a compact strict CY$_3$ threefold $X$
 (projective with $\omega_X \cong \cO_X$ and $h^1(\cO_X) = 0$),
 Chern-character provides a numerical isomorphism
 $K_0^{\mathrm{num}}(X) \otimes \Q \cong \bigoplus_i H^{2i, 2i}(X, \Q)$, so
 \[
 \mathrm{rk}\, K_0^{\mathrm{num}}(X) \;=\; 2 + 2 h^{1,1}(X),
 \]
 accounting for $[\cO_{\mathrm{pt}}]$ and $[X]$ (two) plus one
 divisor/curve pair per Picard generator (two per $h^{1,1}$). K3~$\times$~E
 has $h^{1,1}(K3 \times E) = 21$, hence
 $\mathrm{rk}\, K_0^{\mathrm{num}}(K3 \times E) = 44$. The conifold's rank
 2 is the smallest non-trivial value; the CY$_3$ dimension is $3$ in both
 cases, so rank is independent of $d$.
\end{example}

The stability manifold is stratified by \emph{walls}: codimension-one real
hypersurfaces where the set of semistable objects of a given phase changes.
The open dense complement decomposes into \emph{chambers}, on each of which
the BPS spectrum is locally constant:
\[
 \Stab(\cC) \setminus \{\text{walls}\} \;=\; \bigsqcup_\alpha \mathrm{Ch}_\alpha,
 \qquad \Omega_\alpha \colon \Gamma \to \Z\text{ constant on each } \mathrm{Ch}_\alpha.
\]
A BPS invariant $\Omega_\alpha(\gamma)$ counts (motivically, or virtually)
the stable objects of class $\gamma$ with respect to $\sigma \in \mathrm{Ch}_\alpha$.
Walls are loci where semistable objects of class $\gamma$ develop with
$\Hom$-extensions crossing between chambers.

\begin{definition}[Chamber boundary]
 \label{def:chamber-boundary}
 A \emph{wall of marginal stability} for class $\gamma \in \Gamma$ is the locus
 $W_\gamma = \{\sigma : Z(\gamma) \in \R_{>0} \cdot Z(\gamma')$ for some
 proper sub-object class $\gamma' <_\sigma \gamma\}$. Walls are real
 codimension-one; their collection stratifies $\Stab(\cC)$ into chambers
 $\mathrm{Ch}_\alpha$.
\end{definition}

\subsubsection{The quantum dilogarithm}

Central to the KS wall-crossing formula is the \emph{quantum
dilogarithm}. We derive it from scratch.

\begin{definition}[Quantum dilogarithm]
 \label{def:quantum-dilog}
 For a formal variable $x$ commuting with $\LL^{1/2}$, set
 \[
 \Psi(x) \;\defeq\; \prod_{k \geq 0} \Big(1 - \LL^{k + 1/2}\, x\Big)^{-1}
 \;=\; \sum_{n \geq 0} \frac{\LL^{n^2/2}\, x^n}{\prod_{j=1}^n (1 - \LL^j)}.
 \]
 The second equality is the $q$-binomial theorem / Durfee-square
 identity; we derive it below.
\end{definition}

\begin{proposition}[Durfee-staircase derivation of $\Psi$]
 \label{prop:psi-durfee-square}
 The two formulas for $\Psi(x)$ in \cref{def:quantum-dilog} agree as
 formal power series in $x$ over $\Z[\LL^{\pm 1/2}]$.
\end{proposition}

\begin{proof}[Durfee staircase]
 Expand each factor geometrically:
 \[
 \Psi(x) \;=\; \prod_{k \geq 0} \sum_{n_k \geq 0} \LL^{(k + 1/2) n_k} x^{n_k}
 \;=\; \sum_{(n_0, n_1, \ldots)} \LL^{\sum_k (k + 1/2) n_k}\, x^{\sum_k n_k}.
 \]
 Index the sum by $n \defeq \sum_k n_k$. A sequence $(n_k)_{k \geq 0}$
 with $\sum n_k = n$ parametrises a partition $\lambda = (\lambda_1 \geq
 \lambda_2 \geq \cdots \geq \lambda_n \geq 0)$ into exactly $n$ parts
 $\geq 0$ (allowing $\lambda_i = 0$), where $n_k = \#\{i : \lambda_i =
 k\}$. The exponent of $\LL$ is
 \[
 \sum_k (k + \tfrac{1}{2}) n_k \;=\; \sum_k k\, n_k + \tfrac{n}{2}
 \;=\; |\lambda| + \tfrac{n}{2}.
 \]
 Extract the Durfee staircase $\delta_n = (n-1, n-2, \ldots, 1, 0)$:
 every partition $\lambda$ with $n$ (possibly zero) parts decomposes
 uniquely as $\lambda = \delta_n + \lambda^\bullet$ where $\lambda^\bullet
 = (\lambda_1 - (n-1), \lambda_2 - (n-2), \ldots, \lambda_n - 0)$ is an
 arbitrary partition with $\leq n$ parts $\geq 0$. The staircase
 contributes
 $|\delta_n| = \binom{n}{2} = \frac{n(n-1)}{2}$;
 the residual $\lambda^\bullet$ has generating function
 $\prod_{j=1}^n (1 - \LL^j)^{-1}$ (Euler: partitions into $\leq n$ parts).
 Substituting,
 \[
 \Psi(x)
 \;=\; \sum_{n \geq 0} x^n \LL^{n/2} \cdot \LL^{n(n-1)/2} \cdot
 \prod_{j=1}^n (1 - \LL^j)^{-1}
 \;=\; \sum_{n \geq 0} \frac{\LL^{n^2/2}\, x^n}{\prod_{j=1}^n (1 - \LL^j)},
 \]
 using $\tfrac{n}{2} + \tfrac{n(n-1)}{2} = \tfrac{n^2}{2}$. This is the
 $q$-binomial form of the quantum dilogarithm; the derivation is
 Faddeev--Kashaev (hep-th/9310070) in motivic form, with the staircase
 providing the canonical $\LL^{n^2/2}$ Gaussian shift.
\end{proof}

\begin{remark}[Why the $1/2$-shift]
 The shift $\LL^{k + 1/2}$ in each factor (rather than $\LL^k$ or
 $\LL^{k+1}$) is dictated by the motivic square-root convention and is
 the choice that makes $\Psi$ invariant under the involution
 $\LL \mapsto \LL^{-1}$ up to prefactor. The half-integer exponent
 becomes the Gaussian factor $\LL^{n^2/2}$ in the summation form, which
 is the motivic avatar of the theta-function shadow of the quantum
 dilogarithm at $|\LL| = 1$. In the classical limit $\LL \to 1$,
 $\Psi$ degenerates to $\exp(\mathrm{Li}_2(x))$, the exponentiated
 classical dilogarithm.
\end{remark}

\begin{theorem}[Faddeev--Kashaev pentagon, canonical form]
 \label{thm:psi-pentagon}
 Let $\hat x_0, \hat x_1$ be generators of a rank-two quantum torus with
 relation $\hat x_1 \hat x_0 = \LL\, \hat x_0 \hat x_1$ (equivalently
 $\langle \gamma_1, \gamma_0 \rangle = 1$ in the KS skew-pairing
 convention $\hat x_\gamma \hat x_{\gamma'} = \LL^{\langle \gamma, \gamma'
 \rangle/2} \hat x_{\gamma + \gamma'}$). Then inside $\hat{\mathcal{T}}_\Gamma$,
 \[
 \Psi(\hat x_0)\, \Psi(\hat x_1)
 \;=\; \Psi(\hat x_1)\, \Psi(\LL^{-1/2}\, \hat x_0 \hat x_1)\, \Psi(\hat x_0).
 \]
 The middle factor $\hat x_{\gamma_0 + \gamma_1} = \LL^{-1/2}\, \hat x_0
 \hat x_1$ is the bound-state class $[S_0] + [S_1]$ produced at the wall.
\end{theorem}

\begin{proof}[Derivation from $A_2$ wall-crossing]
 Specialise \cref{gl:thm:ks-wall-crossing} to the $A_2$ sublattice
 $V = \Z \gamma_0 \oplus \Z \gamma_1$ with $\langle \gamma_1, \gamma_0
 \rangle = 1$ (so $\langle \gamma_0, \gamma_1 \rangle = -1$ by
 antisymmetry). Then $\hat x_0 \hat x_1 = \LL^{-1/2}\,
 \hat x_{\gamma_0 + \gamma_1}$ and $\hat x_1 \hat x_0 = \LL^{1/2}\,
 \hat x_{\gamma_0 + \gamma_1}$, so the commutation
 $\hat x_1 \hat x_0 = \LL\, \hat x_0 \hat x_1$ holds and
 $\hat x_{\gamma_0 + \gamma_1} = \LL^{1/2}\, \hat x_1 \hat x_0 =
 \LL^{-1/2}\, \hat x_0 \hat x_1$. Cross a single wall where the phases
 of $Z(\gamma_0), Z(\gamma_1)$ align.
 \begin{itemize}[leftmargin=2em]
 \item LHS chamber: BPS spectrum $\{\gamma_0, \gamma_1\}$ with
 $\Omega = 1$ each; phase order $\arg Z(\gamma_0) <
 \arg Z(\gamma_1)$, producing the increasing-phase ordered product
 $\Psi(\hat x_0) \Psi(\hat x_1)$.
 \item RHS chamber: BPS spectrum
 $\{\gamma_1, \gamma_0 + \gamma_1, \gamma_0\}$ with $\Omega = 1$ on
 each; the bound-state class $[S_0] + [S_1] = \gamma_0 + \gamma_1$
 has materialised; phase order reversed, producing
 $\Psi(\hat x_1) \Psi(\hat x_{\gamma_0 + \gamma_1}) \Psi(\hat x_0)
 = \Psi(\hat x_1) \Psi(\LL^{-1/2} \hat x_0 \hat x_1) \Psi(\hat x_0)$.
 \end{itemize}
 Equality of the two chamber products is the pentagon.

 Algebraic verification (Faddeev--Kashaev hep-th/9310070 (1994)):
 expand both sides as formal series in $\hat x_0, \hat x_1$ using the
 sum form of $\Psi$ and the commutation relation $\hat x_1 \hat x_0 =
 \LL \hat x_0 \hat x_1$; bidegree-by-bidegree matching reduces to the
 $q$-Pascal recursion
 \[
 \binom{m+n}{n}_\LL \;=\; \binom{m+n-1}{n}_\LL + \LL^{m}
 \binom{m+n-1}{n-1}_\LL
 \]
 for the Gaussian binomial coefficients, which is the identity dual to
 the pentagon under the KS motivic Fourier transform. The three BPS
 contributions $\Omega(\gamma_0) + \Omega(\gamma_1) + \Omega(\gamma_0 +
 \gamma_1) = 3$ count the three $\Psi$-factors on the RHS.
\end{proof}

\begin{remark}[Pentagon as BPS conservation law]
 \label{rem:pentagon-bps-conservation}
 The pentagon is the first non-trivial consequence of the KS
 wall-crossing formula and, read as a cluster identity, controls the
 combinatorics of the $A_2$ cluster algebra (Fock--Goncharov,
 Fomin--Zelevinsky). In the DT / motivic dictionary the $A_2$ wall has
 exactly three BPS states with charges $\gamma_1, \gamma_2, \gamma_1 + \gamma_2$;
 the pentagon is the assertion that these three quantum dilogarithms
 re-order consistently. All higher wall-crossings (hexagons, octagons,
 and the continuum of Kontsevich--Soibelman wall crossings for BPS
 rays of irrational slope) are consequences of the pentagon plus the
 cocycle condition \cref{cor:ks-cocycle}.
\end{remark}

\subsubsection{The KS wall-crossing formula}

\begin{theorem}[Kontsevich--Soibelman wall-crossing formula {\cite[\S 6.3]{KS2008}}]
 \label{gl:thm:ks-wall-crossing}
 Let $\sigma_1, \sigma_2 \in \Stab(\cC)$ be points in adjacent chambers
 $\mathrm{Ch}_1, \mathrm{Ch}_2$ separated by a wall $W$. Let $V \subset \Gamma$ be the
 rank-two sublattice of charges whose phases align on $W$, and let
 $\Omega_i(\gamma)$ denote the motivic BPS invariants of chamber $i$.
 Then
 \[
 \prod_{\gamma \in V^+}^{\curvearrowleft} \Psi(\hat{x}_\gamma)^{\Omega_1(\gamma)}
 \;=\;
 \prod_{\gamma \in V^+}^{\curvearrowright} \Psi(\hat{x}_\gamma)^{\Omega_2(\gamma)},
 \]
 where the left (resp.\ right) product is taken in order of decreasing
 (resp.\ increasing) phase of $Z(\gamma)$ at $\sigma_1$ (resp.\ $\sigma_2$).
\end{theorem}

The formula is a chamber-to-chamber gluing relation: the LHS is the product
of quantum-dilogarithm factors built from chamber-$1$ data; the RHS from
chamber-$2$ data; and the formula asserts they are equal as automorphisms of
the completed quantum torus.

\begin{corollary}[Cocycle condition]
 \label{cor:ks-cocycle}
 Passing around a codim-$2$ stratum where three chambers meet along three
 walls, the three pairwise wall-crossing formulas compose to the identity.
 In cohomological language: the quantum-dilogarithm factors $\{\Psi(\hat{x}_\gamma)^{\Omega(\gamma)}\}$
 define a \v{C}ech cocycle for the open cover of $\Stab(\cC)$ by chambers.
 The global BPS spectrum is the cohomology class of this cocycle in
 $\check{H}^1(\Stab(\cC), \Aut(\hat{\mathcal{T}}_\Gamma))$.
\end{corollary}

\begin{proof}[Sketch]
 Crossing three walls and returning to the starting chamber is the identity
 automorphism of $\Stab(\cC)$; the wall-crossing formula applied three
 times must therefore compose to the identity of $\hat{\mathcal{T}}_\Gamma$. This is the
 familiar \v{C}ech 1-cocycle condition on triple overlaps of the
 chamber-indexed cover.
\end{proof}

\begin{remark}[Why this is gluing]
 The data
 \[
 \{\Omega_\alpha(\gamma)\}_{\alpha, \gamma} \longleftrightarrow \{\text{gluing automorphism } \Psi^{\Omega_\alpha}(\gamma) \text{ on wall } W_{\alpha, \alpha+1}\}
 \]
 is the chamber-indexed gluing datum. The chambers are the \emph{local
 charts} of a global BPS sheaf; the quantum-dilogarithm factors are the
 \emph{transition functions}; the wall-crossing formula is the
 \emph{cocycle condition}. This is the same formal structure as
 \v{C}ech-gluing of a line bundle on a manifold, with $U(1)$ replaced by
 $\Aut(\hat{\mathcal{T}}_\Gamma)$ and the sheaf cohomology replaced by pointed
 non-abelian sheaf cohomology $H^1$.
\end{remark}

\subsubsection{The Joyce--Song integration map}

An alternative, equivalent gluing is provided by the Joyce--Song integration
map. We state it at the chain level.

\begin{theorem}[Joyce--Song integration map, Joyce--Song (2012, arXiv:0810.5645) Thm 5.8]
 \label{thm:joyce-song-integration}
 Let $\cC$ be a CY$_3$ category with stability condition $\sigma$ and
 charge lattice $\Gamma$. The moduli stack
 $\mathfrak{M}_\sigma^{\mathrm{ss}}(\gamma)$ of $\sigma$-semistable objects
 of class $\gamma$ is an Artin stack of finite type, and its
 characteristic function $\bar{\delta}_\sigma^{\mathrm{ss}}(\gamma) \in
 \mathrm{SF}^{\mathrm{ind}}(\mathfrak{M})$ is a constructible stack
 function supported on the semistable locus, weighted by the motivic
 Behrend function $\nu_{\mathfrak{M}}$. Joyce--Song (2012, arXiv:0810.5645)
 Thm 5.8 constructs an algebra homomorphism (the \emph{Joyce--Song
 integration map}),
 \[
 I_\sigma \colon \mathrm{SF}^{\mathrm{ind}}_{\mathrm{alg}}(\mathfrak{M})
 \longrightarrow \hat{\mathcal{T}}_\Gamma \otimes_\Z \Q,
 \]
 from the stack-function Hall algebra (with Behrend weighting) to the
 completed quantum torus, characterised on generators by
 \[
 I_\sigma\!\big(\bar{\delta}_\sigma^{\mathrm{ss}}(\gamma)\big)
 \;=\; \bar{\Omega}_\sigma(\gamma)\, \hat{x}_\gamma,
 \qquad
 \bar{\Omega}_\sigma(\gamma) \;=\; \sum_{k \mid \gamma} \frac{1}{k^2}
 \Omega_\sigma(\gamma / k),
 \]
 where $\Omega_\sigma(\gamma) \in \Q$ is the motivic BPS invariant and
 $\bar{\Omega}_\sigma$ is the Joyce--Song ``rational'' DT invariant
 (Joyce--Song Def 5.13). The map $I_\sigma$ is an algebra homomorphism
 (Joyce--Song Thm 5.8): the Hall product
 $\bar{\delta}^{\mathrm{ss}}_\sigma(\gamma_1) \star
 \bar{\delta}^{\mathrm{ss}}_\sigma(\gamma_2)$ on the source (moduli stack
 of extensions) maps to the quantum-torus product
 $\hat{x}_{\gamma_1} \hat{x}_{\gamma_2} = \LL^{\chi(\gamma_1, \gamma_2)}
 \hat{x}_{\gamma_1 + \gamma_2}$ on the target.
\end{theorem}

Chamber-equivariance upgrades the numerical wall-crossing identity to
the stack-function level:

\begin{corollary}[Chamber-equivariance of $I$ (Joyce--Song \S 6.5, Kontsevich--Soibelman \S 6.3)]
 \label{cor:I-chamber-equivariance}
 Under a wall-crossing $\sigma_1 \to \sigma_2$ across wall $W$,
 the integration maps satisfy
 \[
 I_{\sigma_2} \;=\; \mathrm{Ad}_{U_W} \circ I_{\sigma_1},
 \qquad
 U_W \;=\; \prod_{\gamma \in V^+}^{\arg Z_{\sigma_1}\text{-ordered}} \Psi(\hat{x}_\gamma)^{\bar{\Omega}_{\sigma_1}(\gamma)}
 \;\in\; \Aut(\hat{\mathcal{T}}_\Gamma),
 \qquad
 \bar{\Omega}_{\sigma_1}(\gamma) = \sum_{k \mid \gamma} \tfrac{1}{k^2}\,\Omega_{\sigma_1}(\gamma/k),
 \]
 where $V \subset \Gamma$ is the rank-two sublattice whose phases align
 on $W$ and $\bar{\Omega}_{\sigma}$ is the Joyce--Song rational DT
 invariant (Joyce--Song (2012, arXiv:0810.5645) Def.~5.4, Thm.~5.8).
 Equivalently, after the change of coordinates on $\hat{\mathcal{T}}_\Gamma$
 prescribed by $\mathrm{Ad}_{U_W}$, the two chamber images
 $I_{\sigma_1}(\bar{\delta}^{\mathrm{ss}}_{\sigma_1})$ and
 $I_{\sigma_2}(\bar{\delta}^{\mathrm{ss}}_{\sigma_2})$ agree as
 elements of $\hat{\mathcal{T}}_\Gamma$.
\end{corollary}

The corollary upgrades the numerical wall-crossing identity to a statement at
the level of stack functions, and at the level of the target Lie algebra it
says that the generators $\hat{x}_\gamma$ transform by an inner automorphism of
$\hat{U}(\mathfrak{g}_{\mathrm{BPS}})$.

\subsection{Seiberg-duality mutation as gluing}
\label{sec:seiberg-mutation-gluing}

\subsubsection{Mutation at a vertex}

For $Q$ a quiver with potential $(Q, W)$ and $i \in Q_0$ a vertex with no
loops and no $2$-cycles, the \emph{Keller--Yang mutation} $\mu_i(Q, W)$
produces a new quiver with potential, related to $(Q, W)$ by a derived
equivalence of their cluster categories. Geometrically,
mutation-equivalent quivers correspond to Seiberg-dual gauge theories on the
same underlying CY$_3$, and their Jacobi algebras are Morita-equivalent at
the derived level.

\begin{definition}[Quiver mutation]
 \label{def:quiver-mutation}
 The mutation $\mu_i$ at vertex $i$ reverses all arrows at $i$, composes
 paths through $i$ to introduce new arrows $\alpha \beta$ for each
 incoming $\alpha$ and outgoing $\beta$ at $i$, and removes cancelling
 $2$-cycles. The potential $W$ pulls back with the adjusted arrows,
 producing $\mu_i W$.
\end{definition}

\begin{theorem}[Derived equivalence under mutation, Keller--Yang (2011, arXiv:0906.0761) Thm 6.3]
 \label{thm:keller-yang-mutation-wc}
 The Ginzburg CY$_3$ dg-algebras $\Gamma(Q, W)$ and $\Gamma(\mu_i Q,
 \mu_i W)$ have equivalent derived categories
 \[
 \mathrm{Mut}_i^R \colon \mathcal{D}^b(\Gamma(Q, W)) \xrightarrow{\;\sim\;}
 \mathcal{D}^b(\Gamma(\mu_i Q, \mu_i W)),
 \]
 the equivalence being the right-mutation functor at the cluster tilting
 object canonically attached to the quiver. The equivalence restricts to
 the finite-dimensional derived categories
 $D^b_{\mathrm{fd}}(\Gamma(Q,W)) \xrightarrow{\sim}
 D^b_{\mathrm{fd}}(\Gamma(\mu_i Q, \mu_i W))$ of the Jacobi algebras, and
 commutes with the CY$_3$ Serre functor $[3]$ (Keller--Yang
 (2011, arXiv:0906.0761) Thm 6.3).
\end{theorem}

\begin{remark}[Restatement of the \S\ref{sec:gluing:vdb} master theorem]
 \label{rem:keller-yang-restatement}
 The theorem as stated here is a restatement of the \S\ref{sec:gluing:vdb} master
 theorem (\ref{thm:keller-yang-mutation}), specialised to the
 wall-crossing context where the mutation is indexed by the chamber
 boundary rather than by a choice of tilting object. Both forms rest
 on Keller--Yang Keller--Yang (2011, arXiv:0906.0761) (arXiv:0906.0761), whose content
 is the derived-equivalence half of the mutation statement; the
 cluster-level shadow at $K_0$ (Derksen--Weyman--Zelevinsky quiver-mutation
 identity, arXiv:0812.3781) is a numerical consequence, not the
 primary datum.
\end{remark}

\subsubsection{Mutation induces a bialgebra automorphism}
\label{subsec:mutation-bialg-auto}

The Keller--Yang derived equivalence $\mathrm{Mut}_i^R$ induces, by
functoriality of the Davison--Meinhardt critical-cohomology construction,
a bialgebra automorphism of the CoHA.

\begin{theorem}[Mutation-induced bialgebra automorphism]
 \label{thm:mutation-bialgebra-automorphism}
 The right-mutation functor $\mathrm{Mut}_i^R$ of
 \cref{thm:keller-yang-mutation-wc} induces, via the functoriality of the
 critical-cohomology CoHA, a bialgebra automorphism
 \[
 \mu_i^\cH \colon \cH(Q, W) \xrightarrow{\sim} \cH(\mu_i Q, \mu_i W),
 \]
 intertwining the Hall products, the Drinfeld-new coproducts, and the
 bialgebra pairings (at leading residue). On the quantum-torus image,
 $\mu_i^\cH$ acts on generators by
 \[
 \mu_i^\cH(\hat{x}_{\gamma_j}) \;=\; \hat{x}_{\mathrm{FZ}_i(\gamma_j)},
 \]
 where $\mathrm{FZ}_i$ is the Fomin--Zelevinsky cluster mutation
 Fomin--Zelevinsky (2002, math/0104151) at the $i$-th simple on the
 Grothendieck lattice, defined in terms of the signed adjacency matrix
 $B = (b_{ij})$ of the quiver (with $b_{ij} = \#\{\text{arrows } i \to j\} - \#\{\text{arrows } j \to i\}$)
 by $\mathrm{FZ}_i(\gamma_i) = -\gamma_i$ and
 $\mathrm{FZ}_i(\gamma_j) = \gamma_j + [b_{ji}]_+ \cdot \gamma_i$
 for $j \neq i$, where $[x]_+ = \max(x, 0)$. The $K_0$-level action is
 not the naive Weyl reflection on CY$_3$ (where
 $\chi(\gamma_i, \gamma_i) = 0$, \cref{prop:cy3-serre-antisymmetry};
 cf.~also), but the Fomin--Zelevinsky cluster mutation which
 carries the signed arrow-count information through $B$; the lift to
 the critical CoHA is the Davison--Meinhardt Sym-decomposition pullback
 along the derived equivalence.
\end{theorem}

\begin{proof}[From derived Morita to bialgebra automorphism]
 The Keller--Yang right-mutation functor $\mathrm{Mut}_i^R$ is an
 equivalence of CY$_3$ triangulated categories
 Keller--Yang (2011, arXiv:0906.0761), hence induces an equivalence of
 moduli stacks $\mathrm{Mut}_i^R \colon \mathfrak{M}(Q, W) \xrightarrow{\sim}
 \mathfrak{M}(\mu_i Q, \mu_i W)$ compatible with the Koszul potential
 ($W \mapsto \mu_i W$ matches the potential pullback on the Jacobi
 algebra side). Functoriality of critical cohomology
 $\mathfrak{M} \mapsto H^\bullet(\mathfrak{M}, \phi_W)$ under stack equivalences
 preserving the Koszul potential produces the isomorphism of
 critical-cohomology complexes $\cH(Q, W) \xrightarrow{\sim} \cH(\mu_i Q, \mu_i W)$.

 Hall product compatibility: the Hall product is the critical-cohomology
 pushforward along the moduli stack of extensions
 $\mathrm{Ext}_{\mathfrak{M}}$, which is preserved by the equivalence
 $\mathrm{Mut}_i^R$ because extensions in $D^b(\Jac(Q, W))$ correspond
 under derived equivalence to extensions in $D^b(\Jac(\mu_i Q, \mu_i W))$
 up to the canonical Koszul sign.

 Coproduct compatibility: the Drinfeld-new coproduct
 $\Delta(e^{(a)}(z)) = e^{(a)}(z) \otimes 1 + \phi^{+,(a)}(z) \otimes e^{(a)}(z)$
 pulls back along $\mathrm{Mut}_i^R$ to the Drinfeld-new coproduct for
 the mutated generators because the Cartan currents $\phi^{+,(a)}(z)$
 are Serre-duality-derived data and transform under derived equivalence
 by conjugation $\phi^{+,(a)} \mapsto w \phi^{+,(a)} w^{-1}$ with $w$
 the KY tilting element, preserving the coproduct structure.

 Bialgebra pairing (leading residue): the Mukai-pairing integration
 $\int_{\mathfrak{M}} \mathrm{pr}_1^* \cdot \mathrm{pr}_2^* \cdot \phi_W$ is a
 $\Aut(\mathfrak{M})$-invariant bilinear form; the equivalence
 $\mathrm{Mut}_i^R$ acts by an $\Aut(\mathfrak{M})$-element, hence preserves
 the residue. The Fomin--Zelevinsky cluster mutation $\mathrm{FZ}_i$
 on $K_0(\cC)$ is the projection of $\mu_i^\cH$ to the Grothendieck
 group; the lift to $\cH(\cC)$ provides the bialgebra-automorphism
 enhancement.
\end{proof}

\begin{remark}[Etingof--Kazhdan bialgebra-automorphism identity]
 \label{rem:etingof-kazhdan-bialg-auto}
 The mutation-induced $\mu_i^\cH$ is an instance of the more general
 Etingof--Kazhdan classification of bialgebra automorphisms of a QUE
 bialgebra \cite[\S 5]{EtingofKazhdan1996}: every derived Morita
 equivalence of CY$_3$ quiver-with-potential categories descends to
 a classical bialgebra automorphism of their critical CoHAs, with
 $K_0$-shadow the Fomin--Zelevinsky cluster mutation and BPS-Lie-algebra
 shadow a piecewise-linear transformation on $\mathfrak{g}_{\mathrm{BPS}}$
 (the Weyl reflection when the quiver is Dynkin; the cluster
 transformation otherwise). This is the mechanism by which the mutation
 groupoid (Dynkin-ADE: the Weyl group; CY$_3$ in general: the cluster
 groupoid) acts on $\cH(\cC)$ by bialgebra automorphisms.
\end{remark}

\begin{corollary}[Mutation groupoid in $\Aut_{\mathrm{bialg}}(\cH(\cC))$]
 \label{cor:mutation-in-aut-bialg}
 The mutation groupoid, generated by $\{\mu_i^\cH\}_{i \in Q_0}$, is a
 subgroupoid of the bialgebra-automorphism groupoid
 $\Aut_{\mathrm{bialg}}(\cH(\cC))$ for $\cH(\cC) = \cH(Q, W)$. Its action
 on chambers is compatible (via
 \cref{thm:js-bialgebra-chamber-equivariance}) with the KS wall-crossing
 cocycle: for chambers $\mathrm{Ch}_1, \mathrm{Ch}_2$ related by a sequence of
 mutations $\mu_{i_1} \cdots \mu_{i_r}$,
 \[
 \mu_{i_r}^\cH \circ \cdots \circ \mu_{i_1}^\cH
 \;=\; \mathrm{Ad}_{U_{W_{12}}},
 \]
 where $U_{W_{12}} = \prod_j \Psi(\hat{x}_{\gamma_j})^{\Omega_{12}(\gamma_j)}$
 is the KS ordered product across the intermediate walls between
 $\mathrm{Ch}_1$ and $\mathrm{Ch}_2$.
\end{corollary}

\subsubsection{Conifold Seiberg-duality chain: NCDT $\to$ DT $\to$ PT}
\label{subsec:conifold-seiberg-chain}

For the resolved conifold $\bar{Y}$ the stability manifold is
four-real-dimensional with three canonical chambers related by
Seiberg-duality mutations at the two vertices of the Klebanov--Witten
quiver:

\begin{theorem}[Conifold three-chamber chain]
 \label{thm:conifold-seiberg-chain}
 emph{(Nagao--Nakajima Nagao--Nakajima (2011, arXiv:0809.2992).)}
 The stability manifold $\Stab(D^b(\bar{Y}))$ contains three adjacent
 chambers
 \[
 \mathrm{Ch}_0 \;\xleftrightarrow{\mu_0}\; \mathrm{Ch}_1 \;\xleftrightarrow{\mu_1}\; \mathrm{Ch}_2
 \]
 with BPS spectra
 \begin{itemize}[leftmargin=2em]
 \item $\mathrm{Ch}_0$: NCDT (non-commutative Donaldson--Thomas),
 $\Omega_0(\gamma_0) = \Omega_0(\gamma_1) = 1$,
 $\Omega_0(k\gamma_0 + \ell\gamma_1) = -1$ for $k, \ell \geq 1$;
 \item $\mathrm{Ch}_1$: DT (Donaldson--Thomas),
 $\Omega_1(\gamma_0) = \Omega_1(\gamma_1) = 1$,
 $\Omega_1(k(\gamma_0 + \gamma_1)) = -2$ for $k \geq 1$,
 generating function $Z^{\mathrm{DT}} = M(-q)^2 \prod_k (1 - Q(-q)^k)^k$;
 \item $\mathrm{Ch}_2$: PT (Pandharipande--Thomas),
 $\Omega_2(\gamma_0) = \Omega_2(\gamma_1) = 1$,
 $\Omega_2(k(\gamma_0 + \gamma_1)) = -2$ for $k \geq 1$ (same
 curve BPS count), but phase-ordered differently;
 generating function $Z^{\mathrm{PT}} = \prod_k (1 - Q(-q)^k)^k$.
 \end{itemize}
 The mutations $\mu_0, \mu_1$ at the two vertices of the KW quiver induce
 the bialgebra automorphisms $\mu_0^\cH, \mu_1^\cH \in
 \Aut_{\mathrm{bialg}}(\cH(\bar{Y}))$ realising the three chamber transitions.
\end{theorem}

\begin{proof}[Chain-level]
 Chamber structure is Nagao--Nakajima Nagao--Nakajima (2011, arXiv:0809.2992)
 and Bridgeland Bridgeland (2011, arXiv:1002.4372) arXiv:1002.4372; NCDT is
 Szendr\H{o}i Szendr\H{o}i (2008, arXiv:0705.3419) arXiv:0705.3419. The DT/PT transform
 $Z^{\mathrm{DT}}/Z^{\mathrm{PT}} = M(-q)^2$ is the Behrend-weighted generating
 function of $0$-dimensional subschemes of $\bar{Y}$
 (Bryan--Morrison (2009, arXiv:0812.5007) Thm.~1.1). The
 induced bialgebra automorphisms are the
 \cref{thm:mutation-bialgebra-automorphism} instantiation with
 $i \in \{0, 1\}$ the KW vertices.

 Explicitly, $\mu_0$ reverses the arrows $a_{1,2}\colon 0 \to 1$ and
 $b_{1,2}\colon 1 \to 0$ at vertex $0$, producing the mutated quiver
 with arrows $a_{1,2}^*\colon 1 \to 0$ and $b_{1,2}^*\colon 0 \to 1$
 plus introduced composites; the KW potential
 $W_{\mathrm{con}} = a_1 b_1 a_2 b_2 - a_1 b_2 a_2 b_1$ pulls back to
 $\mu_0 W_{\mathrm{con}}$ via the Keller--Yang prescription.

 At the Grothendieck-lattice level, the KW $B$-matrix
 $b_{ij} = \#\{\text{arrows } i \to j\} - \#\{\text{arrows } j \to i\} = 2 - 2 = 0$
 is symmetric (both directions carry two arrows); FZ mutation
 on a symmetric $B$-matrix fixes the simple roots at the Grothendieck-lattice
 level: $\mathrm{FZ}_0(\gamma_j) = \gamma_j$ for $j \neq 0$ and
 $\mathrm{FZ}_0(\gamma_0) = -\gamma_0$ only on the mutated vertex. The
 non-trivial content of the mutation lives \emph{above} $K_0$: at the
 critical-cohomology level, $\mu_0^\cH$ exchanges the NCDT and DT
 Sym-decompositions of $\cH(\bar Y)$ by permuting the BPS
 tensor-factors along the Fomin--Zelevinsky cluster exchange. The
 $K_0$-shadow is the identity; the bialgebra-automorphism is
 non-trivial on $\cH(\bar Y)$.

 Compatibility with chamber-transported pairings
 (\cref{thm:hopf-pairing-chamber-invariance}): the Cartan currents
 $H(z) = \phi^{+,(0)}(z)$, $K(z) = \phi^{+,(1)}(z)$ are exchanged
 under the conifold $\Z/2$-symmetry $\tau$ swapping the two vertices;
 mutation $\mu_0$ fixes $K$ and acts on $H$ by the cluster
 transformation $H(z) \mapsto H(z)^{-1}$ (the FZ inverse on the
 single mutated coordinate), preserving the leading residue of the
 pairing. The all-mode preservation remains conjectural at the full-mode level.
\end{proof}

\begin{remark}[Pure mutations on the conifold generate $B_2 \cong \Z$, not $B_3$]
 \label{rem:conifold-pure-mutations-B2}
 The Klebanov--Witten quiver has $|Q_0| = 2$ vertices. The two
 mutations $\mu_0^\cH, \mu_1^\cH \in \Aut_{\mathrm{bialg}}(\cH(\bar Y))$
 are involutions on $K_0$, and their cluster-algebra interaction is
 the rank-two $A_1^{(1)}$-affine story: the product $\mu_0 \mu_1$ has
 infinite order. The subgroup
 $\langle \mu_0^\cH, \mu_1^\cH \rangle \subseteq \Aut_{\mathrm{bialg}}(\cH(\bar Y))$
 generated by pure mutations is cyclic,
 $B_{|Q_0|} = B_2 \cong \Z$, on the single generator $\mu_0 \mu_1$
 (Nagao--Nakajima arXiv:0809.2992). A braid relation
 $\sigma_1 \sigma_2 \sigma_1 = \sigma_2 \sigma_1 \sigma_2$ defining
 $B_3$ requires three strands and cannot arise from two mutation
 generators alone.
\end{remark}

\begin{conjecture}[$B_3 \subseteq \Aut_{\mathrm{bialg}}(\cH(\bar{Y}))$ via mutations and taste shift]
 \label{conj:braid-b3-conifold}
 The bialgebra-automorphism group $\Aut_{\mathrm{bialg}}(\cH(\bar{Y}))$
 contains a subgroup isomorphic to $B_3$, generated by the three
 elements $\{\mu_0^\cH, \mu_1^\cH, \tau^\cH\}$, where $\tau^\cH$ is
 the \emph{taste-shift} automorphism induced by the
 $\Z$-shift on $D^b(\bar Y)$ along the
 $\Aut^0(\bar Y) = \mathbb{C}^\times$-rotation of the $\mathbb{P}^1$
 zero-section (equivalently, the cyclic generator of the Picard
 lattice on the $\mathcal{O}_{\mathbb{P}^1}(-1)$ factor). Setting
 $\sigma_1 = \mu_0^\cH \tau^\cH$ and $\sigma_2 = \mu_1^\cH \tau^\cH$,
 the braid relation
 \[
 \sigma_1 \sigma_2 \sigma_1 \;=\; \sigma_2 \sigma_1 \sigma_2
 \]
 is conjectured to hold in $\Aut_{\mathrm{bialg}}(\cH(\bar Y))$,
 realising a $B_3$-action. The pure-mutation subgroup
 $\langle \mu_0^\cH, \mu_1^\cH\rangle$ is only $B_2 \cong \Z$
 (\cref{rem:conifold-pure-mutations-B2}); the taste generator
 $\tau^\cH$ supplies the third strand.
\end{conjecture}

\begin{remark}[Primary source and scope]
 \label{rem:B3-conifold-primary-source}
 The conifold quiver is rank-two $A_1^{(1)}$-affine; the $K_0$-level
 cluster combinatorics is Fomin--Zelevinsky (2002, math/0104151). The
 lift to a genuine $B_3$-action on $\cH(\bar Y)$ at the
 critical-cohomology level is open: no published theorem asserts it,
 and the taste generator $\tau^\cH$ is the geometric candidate third
 strand derived from $\Aut^0(\bar Y) = \mathbb{C}^\times$.
 Automorphism-compatibility is stated in $\Aut_{\mathrm{bialg}}$
 because $\cH(\bar Y)$ is a bialgebra; antipode compatibility is a
 statement on the double
 $D(\cH(\bar Y)) = Y_\hbar(\widehat{\mathfrak{gl}}(1|1))$.
\end{remark}

\subsubsection{Weyl group action on chambers}

For a simply-laced Dynkin quiver $Q$, iterated mutations generate the Weyl
group $W(Q)$. More generally, for quivers arising from CY$_3$ via tilting,
mutations generate a \emph{mutation groupoid} acting on the space of
tilting objects; this groupoid acts on $\Stab(\cC)$ by permuting chambers.
Explicitly:

\begin{proposition}[Mutation action on chambers]
 \label{prop:mutation-on-chambers}
 The mutation $\mu_i$ induces a homeomorphism $\mu_i \colon \Stab(\cC) \to \Stab(\cC)$
 permuting chambers; two chambers $\mathrm{Ch}_1$, $\mathrm{Ch}_2$ lie in the same
 $\mu$-orbit if and only if the tilting objects $T_1$, $T_2$ realising them
 are related by iterated mutation.
\end{proposition}

\begin{corollary}[Weyl action on chamber lattice]
 \label{cor:weyl-chamber-action}
 For $Q$ of ADE type with acyclic orientation, the mutation groupoid is
 exactly $W(Q)$. The action on the chamber lattice matches the reflection
 action of $W(Q)$ on the root lattice of the associated simple Lie algebra.
\end{corollary}

\subsubsection{Chamber-independence of the CoHA}

The key gluing theorem connects the local (chamber-by-chamber) and
global (chamber-free) pictures:

\begin{theorem}[Chamber-independence of CoHA, Davison--Meinhardt (2020, arXiv:1601.02479) Thm A+B]
 \label{thm:coha-chamber-independent}
 Let $\cH(\cC)$ denote the Davison--Meinhardt critical CoHA of a CY$_3$
 category (Thm A). Then $\cH(\cC)$ is canonically independent of the
 stability condition $\sigma \in \Stab(\cC)$. The Sym-decomposition
 (Thm B / integrality) is $\sigma$-dependent, but the underlying bialgebra
 is not: different chambers give \emph{isomorphic} Hall algebras with the
 isomorphism $\phi_{12} \colon \cH_{\sigma_1}(\cC) \xrightarrow{\sim}
 \cH_{\sigma_2}(\cC)$ realised explicitly by the adjoint action of the
 KS wall-crossing cocycle.
\end{theorem}

\begin{proof}[Chain-level sketch via Davison--Meinhardt Thm A+B]
 The Davison--Meinhardt (2020, arXiv:1601.02479) critical CoHA is
 $\cH(\cC) \defeq H^\bullet(\mathfrak{M}(\cC), \phi_W)$,
 the critical cohomology of the moduli stack of objects with its
 Koszul potential $W$ induced by the CY$_3$ pairing. This datum is
 intrinsic to the dg-structure of $\cC$ and does not reference any
 stability condition $\sigma$: Davison--Meinhardt Thm A constructs
 $\cH(\cC)$ directly from $(\cC, W)$ without choosing $\sigma$. What
 $\sigma$ furnishes is a \emph{Sym-decomposition} (Davison--Meinhardt
 Thm B / integrality),
 \[
 \cH(\cC) \;\cong\; \bigotimes_{\gamma \in \Gamma^+}^{\arg Z_\sigma(\gamma)\text{-ordered}}
 \Sym^\bullet\!\Big(\mathrm{BPS}_\sigma(\gamma) \otimes H^\bullet(\B \G_m)\Big),
 \]
 where $\mathrm{BPS}_\sigma(\gamma)$ is the BPS cohomology in charge
 $\gamma$ at stability $\sigma$ and the tensor factor $H^\bullet(\B \G_m) =
 \Q[u]$ records the overall $\G_m$-equivariant direction. Two chambers
 $\sigma_1, \sigma_2$ produce distinct $\sigma$-ordered Sym-decompositions
 of the \emph{same} bialgebra $\cH(\cC)$; the KS wall-crossing formula is
 the algebra isomorphism $\phi_{12} \colon \cH_{\sigma_1}(\cC)
 \xrightarrow{\sim} \cH_{\sigma_2}(\cC)$ conjugating one decomposition
 into the other, realised at the motivic-cohomology level by the adjoint
 action of the ordered product $\prod_\gamma \Psi(\hat{x}_\gamma)^{\Omega_1(\gamma)}$.
 The underlying algebra is the same; only its Sym-presentation changes.
\end{proof}

\begin{corollary}[Representations are chamber-parametrised]
 \label{cor:reps-chamber-parametrised}
 While $\cH(\cC)$ is chamber-independent, individual irreducible
 representations of $\cH(\cC)$ depend on the chamber. A BPS brick
 $\mathrm{BPS}_\sigma(\gamma)$ is an irreducible $\cH(\cC)$-module in
 chamber $\sigma$; its $\cH(\cC)$-module structure changes under
 wall-crossing by the KS cocycle.
\end{corollary}

\subsubsection{Shifted quiver Yangians index chambers}

Galakhov--Li--Yamazaki (GLY) Galakhov--Li--Yamazaki (2021, arXiv:2008.07006)
Galakhov--Li--Yamazaki (2022, arXiv:2108.10286) constructed the \emph{shifted quiver Yangian}
$Y_{\hbar, \zeta}(Q, W)$ associated to a quiver-with-potential $(Q, W)$ and
a \emph{shift parameter} $\zeta \in \Gamma^\vee$. The shift parameter
precisely records the chamber data:

\begin{theorem}[Shifted quiver Yangian indexes chambers, GLY]
 \label{thm:shifted-yangian-chambers}
 Let $\sigma \in \Stab(\cC)$ lie in chamber $\mathrm{Ch}_\sigma$. Associated to
 $\sigma$ is a shift $\zeta_\sigma \in \Gamma^\vee$, and the corresponding
 representation-theoretic CoHA-image is the shifted quiver Yangian
 $Y_{\hbar, \zeta_\sigma}(Q, W)$. Wall-crossing changes $\zeta$ by the
 KS-$\Psi$-shift; the abstract Yangian
 $Y_{\hbar, \zeta}(Q, W)$ at generic $\zeta$ is chamber-independent, but
 the \emph{representation at a specific BPS spectrum} is
 $\zeta_\sigma$-dependent.
\end{theorem}

The shift parameter is thus a \emph{local coordinate on the chamber lattice};
the shifted Yangian provides a concrete parametrised algebra whose specialisation
at $\zeta_\sigma$ realises the BPS-brick module of chamber $\sigma$.

\subsection{Joyce--Song integration and chamber-invariance of the bialgebra pairing}
\label{sec:js-hopf-chamber-independence}

\subsubsection{The positive-half bialgebra and its Drinfeld double}
\label{subsec:bialgebra-not-hopf}

The critical CoHA $\cH(\cC)$ on a CY$_3$ category is a $\Z_2$-graded
\emph{associative bialgebra} in super-vector spaces over $\Q((\hbar))$. Its
four structure maps
\[
 \mu \colon \cH(\cC) \otimes \cH(\cC) \to \cH(\cC),
 \quad
 \iota \colon \Q \to \cH(\cC),
 \quad
 \Delta \colon \cH(\cC) \to \cH(\cC) \otimes \cH(\cC),
 \quad
 \epsilon \colon \cH(\cC) \to \Q,
\]
are the Davison--Meinhardt Hall product, the vacuum inclusion, the
Drinfeld-new coproduct
$\Delta(e^{(a)}(z)) = e^{(a)}(z) \otimes 1 + \phi^{+,(a)}(z) \otimes e^{(a)}(z)$,
and the standard augmentation. An antipode $S$ on $\cH(\cC)$ itself is
obstructed: $\cH(\cC)$ is the positive half of a larger Hopf-algebraic
completion and satisfies the bialgebra axioms without the Hopf axiom.
The antipode is installed only after passage to the Drinfeld double
\[
 D(\cH(\cC)) \;=\; \cH(\cC) \otimes_{\cH^0} \cH(\cC)^{\mathrm{op, cop}},
\]
paired with the negative-half shadow $\cH(\cC)^{\mathrm{op, cop}}$ along
the bialgebra pairing below. Etingof--Kazhdan
\cite[Prop.~3.4]{EtingofKazhdan1996} derive the double's antipode from
a non-degenerate bialgebra pairing: $S$ is the unique $\Q((\hbar))$-linear
map on $D(\cH)$ intertwining the pairing of $\cH$ with
$\cH^{\mathrm{op,cop}}$, characterised by
$\mu \circ (S \otimes \id) \circ \Delta = \iota \circ \epsilon$.

The bilinear form the chamber-invariance statement tracks is therefore
a \emph{bialgebra pairing}
\[
 (-, -)_\cH \colon \cH(\cC) \otimes \cH(\cC) \longrightarrow \Q((\hbar)),
\]
in the Etingof--Kazhdan sense \cite[\S 3]{EtingofKazhdan1996}: bilinear,
compatible with $\mu$ and $\Delta$ by
$(\xi \eta, \zeta)_\cH = (\xi \otimes \eta, \Delta \zeta)_\cH$, and
non-degenerate on each $K_0$-graded piece. Antipode-compatibility
$(S\xi, \eta)_\cH = (\xi, S\eta)_\cH$ is a statement on the double, not
on $\cH(\cC)$. For elements $[E], [F] \in \cH(\cC)$ represented by
critical cohomology classes on the moduli stack,
\[
 ([E], [F])_\cH \;=\; \int_{\mathfrak{M}} \mathrm{pr}_1^*[E] \cdot \mathrm{pr}_2^*[F] \cdot \phi_W,
\]
where $\phi_W$ is the vanishing-cycle functor for the Koszul potential
$W$ and the integral is the equivariant pushforward to the point. On the
conifold, $D(\cH(\bar{Y})) = Y_{\hbar}(\widehat{\mathfrak{gl}}(1|1))$ is
the full quantum super-Yangian with antipode
$S(e^{(a)}(z)) = -\phi^{+,(a)}(z)^{-1} e^{(a)}(z) \phi^{+,(a)}(z)$; the
positive half
$\cH(\bar Y) = Y^+_{\hbar}(\widehat{\mathfrak{gl}}(1|1))^{\mathrm{con}}$
carries only $(\mu, \iota, \Delta, \epsilon)$.

\begin{remark}[Hopf structure stratified by equivariance]
 \label{rem:hopf-scope-stratified}
 The double $D(\cH(\cC))$ is strict $\Z_2$-graded Hopf at the rational,
 trigonometric, and toroidal equivariance strata, by the
 Etingof--Kazhdan quantisation theorem
 \cite{EtingofKazhdan1996}. At the elliptic stratum the
 double is quasi-Hopf, carrying the Felder--Jimbo--Konno dynamical
 associator: the coassociator $\Phi(\lambda) \in D(\cH(\cC))^{\otimes 3}$
 depends on a dynamical Cartan parameter $\lambda \in \mathfrak{h}^*$
 and satisfies the pentagon up to the Weyl-translation shift
 $\lambda \mapsto \lambda + h^{(i)}$ on the $i$-th tensor factor (Felder's
 dynamical elliptic $R$-matrix, Felder (1994, hep-th/9407154); elliptic
 quantum affine instantiation, Jimbo--Konno (1997, q-alg/9610013)).
 Chamber-invariance in this section concerns the bialgebra $\cH(\cC)$
 and $(-,-)_\cH$; the stratum-dependent (quasi-)Hopf structure of
 $D(\cH(\cC))$ is its downstream consequence.
\end{remark}

\subsubsection{Chamber-invariance of the bialgebra pairing}
\label{subsec:chamber-invariance-pairing}

Decomposing along the $K_0$-grading, the bialgebra pairing splits into
homogeneous residues
\[
 c_{ab}(m, n) \;=\; (e^{(a)}_m, f^{(b)}_n)_\cH,
 \qquad
 m, n \in \Z, \ \ (a, b) \in I \times I,
\]
where $m+n$ tracks the $K_0$-degree sum. Chamber-invariance stratifies:
the leading residue ($m+n=0$, $K_0$-orthogonality regime) is preserved
by every wall-crossing automorphism as a theorem; propagation to all
modes is conjectural.

\begin{theorem}[Residue-level chamber-invariance of $(-, -)_\cH$, Etingof--Kazhdan]
 \label{thm:hopf-pairing-chamber-invariance}
 Let $u = U_W \in \hat{\mathcal{T}}_\Gamma \subset \cH(\cC)$ be the wall-crossing
 element of \cref{thm:coha-chamber-independent}. On each leading-residue
 matrix element $c_{ab}(m, n)$ with $m + n = 0$ (i.e.\ on the $K_0$-diagonal),
 \[
 \big(\mathrm{Ad}_u(\xi_{m}^{(a)}),\, \mathrm{Ad}_u(\eta_{n}^{(b)})\big)_\cH
 \;=\; (\xi_{m}^{(a)},\, \eta_{n}^{(b)})_\cH
 \qquad \text{whenever } m + n = 0 \text{ in } K_0.
 \]
 Consequently the leading-residue coefficients $c_{ab}(0, 0)$ are
 chamber-free: they depend only on the bialgebra $\cH(\cC)$, not on
 $\sigma \in \Stab(\cC)$.
\end{theorem}

\begin{proof}[$K_0$-orthogonality of $\alpha_W$ corrections]
 The bialgebra pairing is $K_0$-diagonal:
 $(\xi_\gamma, \eta_{\gamma'})_\cH = 0$ unless $\gamma + \gamma' = 0$
 (Serre-duality complementary charges,
 \cref{prop:cy3-serre-antisymmetry}). Expand the wall-crossing
 automorphism as $\mathrm{Ad}_u = \id + \alpha_W$ with correction
 $\alpha_W$ landing strictly inside the positive cone
 $V^+ = \Z_{\geq 0}\gamma_1 \oplus \Z_{\geq 0}\gamma_2$ of the rank-two
 sublattice $V$ through $W$:
 \[
 \alpha_W(\xi_\alpha) \;\in\; \bigoplus_{\substack{k, \ell \geq 0 \\ k + \ell \geq 1}}
 \cH(\cC)_{\alpha + k\gamma_1 + \ell\gamma_2}.
 \]
 For $\xi_\alpha, \eta_\beta$ with $\alpha + \beta = 0$, bilinearity
 gives
 \[
 (\mathrm{Ad}_u \xi_\alpha, \mathrm{Ad}_u \eta_\beta)_\cH
 = (\xi_\alpha, \eta_\beta)_\cH + (\alpha_W \xi_\alpha, \eta_\beta)_\cH
 + (\xi_\alpha, \alpha_W \eta_\beta)_\cH
 + (\alpha_W \xi_\alpha, \alpha_W \eta_\beta)_\cH.
 \]
 The three correction terms pair elements of $K_0$-degrees
 $(\alpha + k\gamma_1 + \ell\gamma_2, \beta)$,
 $(\alpha, \beta + k'\gamma_1 + \ell'\gamma_2)$, and
 $(\alpha + k\gamma_1 + \ell\gamma_2, \beta + k'\gamma_1 + \ell'\gamma_2)$
 with $(k, \ell) \neq (0, 0)$ or $(k', \ell') \neq (0, 0)$; the total
 degree sum is strictly positive in $V^+$, so $K_0$-orthogonality
 forces each correction term to vanish. Thus
 $(\mathrm{Ad}_u \xi_\alpha, \mathrm{Ad}_u \eta_\beta)_\cH
 = (\xi_\alpha, \eta_\beta)_\cH$ on the $K_0$-diagonal.
 The identity transports from $\hat{\mathcal{T}}_\Gamma$ to $\cH(\cC)$ along
 the Davison--Meinhardt critical-cohomology lift: the Sym-decomposition
 is natural in $\sigma$ and the projection
 $\cH(\cC) \twoheadrightarrow \hat{\mathcal{T}}_\Gamma$ commutes with
 $\mathrm{Ad}_u$ on the $K_0$-diagonal. Etingof--Kazhdan
 \cite[Prop.~3.4]{EtingofKazhdan1996} supplies the classical-limit
 $K_0$-orthogonality; Davison--Meinhardt (2020, arXiv:1601.02479)
 supplies the critical-cohomology instantiation.
\end{proof}

\begin{conjecture}[All-mode chamber-invariance, Etingof--Kazhdan + GLY]
 \label{conj:hopf-pairing-all-mode-chamber-invariance}
 emph{(S5-level conjecture.)}
 The full bialgebra pairing is chamber-invariant at all mode indices:
 for every $\xi, \eta \in \cH(\cC)$,
 \[
 \big(\mathrm{Ad}_{U_W}(\xi),\, \mathrm{Ad}_{U_W}(\eta)\big)_\cH
 \;=\; (\xi,\, \eta)_\cH.
 \]
 Equivalently, the asymmetric super-trace correction
 $c_{ab}(m, n) \mapsto c_{ab}(m, n) \cdot (1 + \hbar \cdot c_{ab} \cdot \delta_{m + n, 0})$
 with $c_{00} = 0$, $c_{11} = -1$ on $\widehat{\mathfrak{gl}}(1|1)^{\mathrm{con}}$
 is preserved under wall-crossing at every mode, including the leading
 residue.
\end{conjecture}

\begin{remark}[Why the all-mode statement is conjectural]
 \label{rem:all-mode-conjectural}
 The $c_{11} = -1$ asymmetric correction on the odd root's Killing
 form is derived from the super-trace sign in the Davison super-shuffle
 presentation and is consistent with $\widehat{\mathfrak{gl}}(1|1)$
 structure; it has \emph{not} been verified against published
 $Y(\mathfrak{gl}(m|n))$ Hopf pairings (Nazarov 1991,
 Arnaudon--Cramp\'e--Doikou--Frappat--Ragoucy 2003, Gow 2005, Peng 2011).
 The wall-crossing-compatibility of the all-mode propagated correction
 is therefore an S5 conjecture pending primary-source verification.
 The residue-level \cref{thm:hopf-pairing-chamber-invariance} is
 established independently via $K_0$-orthogonality of the
 wall-crossing corrections.
\end{remark}

\begin{proof}[Geometric proof of \cref{thm:hopf-pairing-chamber-invariance} via Behrend-weighted symmetry]
 A second derivation of the residue-level statement: the wall-crossing
 automorphism $\phi_{12} = \mathrm{Ad}_{U_W}$ is inner in
 $\Aut_{\mathrm{bialg}}(\cH(\cC))$, with
 $U_W = \prod_\gamma \Psi(\hat{x}_\gamma)^{\bar{\Omega}_{\sigma_1}(\gamma)}
 \in \hat{\mathcal{T}}_\Gamma$ and $\bar{\Omega}_\sigma$ the rational DT
 invariants (Joyce--Song (2012, arXiv:0810.5645) \S 5.2). The
 Behrend-weighted integration
 $\int_{\mathfrak{M}} \mathrm{pr}_1^*(-) \cdot \mathrm{pr}_2^*(-) \cdot \phi_W$
 computing $(-, -)_\cH$ is $\Aut(\mathfrak{M})$-invariant under symmetries
 of the moduli stack, and the CY$_3$ property of $\phi_W$ delivers the
 symmetry required: the vanishing-cycle sheaf $\phi_W$ for a CY$_3$
 Koszul potential carries a canonical self-duality (Brav--Bussi--Dupont--
 Joyce--Szendr\H{o}i, arXiv:1211.3259), which combined with $K_0$-grading
 preservation forces $(\mathrm{Ad}_u \xi, \mathrm{Ad}_u \eta)_\cH
 = (\xi, \eta)_\cH$ on $K_0$-complementary charges. The argument extends
 to higher modes under
 \cref{conj:hopf-pairing-all-mode-chamber-invariance}.
\end{proof}

\begin{corollary}[Bialgebra Hilbert series is chamber-independent]
 \label{cor:hilbert-series-chamber-indep}
 The Hilbert series $\sum_\gamma \dim \cH^\gamma(\cC) \cdot \hat{x}_\gamma$
 has chamber-independent coefficients, and the generating function of
 leading-residue pairing values $([E], [E])_\cH|_{\mathrm{res}}$ on any
 fixed basis is chamber-invariant. The all-mode generating function is
 chamber-invariant conditional on
 \cref{conj:hopf-pairing-all-mode-chamber-invariance}.
\end{corollary}

\subsubsection{Joyce--Song integration map is chamber-equivariant}
\label{subsec:js-integration-chamber-equivariance}

The Joyce--Song integration map \cref{thm:joyce-song-integration} upgrades
the numerical wall-crossing relation \cref{gl:thm:ks-wall-crossing} to a
statement at the level of stack functions; at the bialgebra level it is
the chamber-equivariance statement:

\begin{theorem}[Joyce--Song Thm 5.8 chamber-equivariance, bialgebra form]
 \label{thm:js-bialgebra-chamber-equivariance}
 For any wall-crossing $\sigma_1 \to \sigma_2$ across wall $W$,
 \[
 I_{\sigma_2} \;=\; \mathrm{Ad}_{U_W} \circ I_{\sigma_1},
 \qquad
 U_W \;=\; \prod_{\gamma \in V^+}^{\arg Z_{\sigma_1}\text{-ordered}}
 \Psi(\hat{x}_\gamma)^{\bar{\Omega}_{\sigma_1}(\gamma)},
 \qquad
 \bar{\Omega}_{\sigma_1}(\gamma) \;=\; \sum_{k \mid \gamma} \tfrac{1}{k^2}\,\Omega_{\sigma_1}(\gamma/k),
 \]
 as algebra homomorphisms from the stack-function Hall algebra to the
 completed quantum torus; the exponent is the Joyce--Song rational DT
 invariant, not the integer $\Omega$. Three assertions compose:
 \begin{enumerate}[leftmargin=2em]
 \item\label{js1} \emph{(Joyce--Song (2012, arXiv:0810.5645) Thm 5.8)}
 $I_\sigma$ is an algebra homomorphism for every $\sigma$, and the
 identity $I_{\sigma_2} = \mathrm{Ad}_{U_W} \circ I_{\sigma_1}$
 holds on generators.
 \item\label{js2} \emph{(Kontsevich--Soibelman \cite[\S 6.3]{KS2008})}
 $U_W \in \Aut_{\mathrm{bialg}}(\hat{\mathcal{T}}_\Gamma)$ is a bialgebra
 automorphism (not merely an algebra automorphism) of the
 completed quantum torus.
 \item\label{js3} \emph{(Davison--Meinhardt (2020, arXiv:1601.02479) Thm A+B)}
 The commutative triangle of \ref{js1} and \ref{js2} lifts from
 the $K_0$-graded quotient $\hat{\mathcal{T}}_\Gamma$ to the full critical CoHA
 $\cH(\cC)$ via the Sym-decomposition, preserving the bialgebra
 structure pointwise in each $K_0$-degree.
 \end{enumerate}
 The composite is the statement that the ordered product
 $U_W = \prod \Psi(\hat{x}_\gamma)^{\bar{\Omega}_{\sigma_1}(\gamma)}$ is
 the explicit bialgebra-automorphism witness of the chamber change
 $\sigma_1 \to \sigma_2$; the rational-DT exponent $\bar{\Omega}$
 absorbs the multicover strict-semistability corrections
 (Joyce--Song (2012, arXiv:0810.5645) \S 5.2, Def.~5.4).
\end{theorem}

\begin{proof}[Structural]
 \ref{js1} is Joyce--Song Joyce--Song (2012, arXiv:0810.5645): the Hall product
 $\bar{\delta}_\sigma^{\mathrm{ss}}(\gamma_1) \star \bar{\delta}_\sigma^{\mathrm{ss}}(\gamma_2)$
 on the stack-function side maps to
 $\bar{\Omega}_\sigma(\gamma_1) \bar{\Omega}_\sigma(\gamma_2) \hat{x}_{\gamma_1} \hat{x}_{\gamma_2}
 = \LL^{\chi(\gamma_1, \gamma_2)} \bar{\Omega}_\sigma(\gamma_1) \bar{\Omega}_\sigma(\gamma_2) \hat{x}_{\gamma_1 + \gamma_2}$
 on the quantum-torus side, compatibly with the Behrend-weighted integration.

 \ref{js2} is the observation that each $\Psi(\hat{x}_\gamma)$ is a bialgebra
 automorphism of $\hat{\mathcal{T}}_\Gamma$: it preserves the coproduct
 $\Delta(\hat{x}_\gamma) = \hat{x}_\gamma \otimes \hat{x}_\gamma$ (grouplike) and
 intertwines the Euler-form-twisted multiplication, a direct
 verification on generators using the pentagon identity
 \cref{thm:psi-pentagon}. Products of bialgebra automorphisms are
 bialgebra automorphisms; hence $U_W \in \Aut_{\mathrm{bialg}}(\hat{\mathcal{T}}_\Gamma)$.

 \ref{js3} is the critical-cohomology upgrade. Davison--Meinhardt
 Davison--Meinhardt (2020, arXiv:1601.02479) show that the Sym-decomposition
 $\cH(\cC) \cong \bigotimes \Sym^\bullet(\mathrm{BPS}_\sigma \otimes H^\bullet(\B\G_m))$
 is natural in $\sigma$: different chambers permute the tensor factors,
 and the chamber-isomorphism $\phi_{12}$ of
 \cref{thm:coha-chamber-independent} intertwines the two
 Sym-decompositions. The induced map on $K_0$-quotients is exactly
 $\mathrm{Ad}_{U_W}$; the lift to $\cH(\cC)$ preserves the bialgebra
 structure because each Sym factor is a commutative Hopf ideal in the
 critical-cohomology complex (Davison--Meinhardt integrality).

 The three assertions compose to the claim that $U_W$ is the explicit
 bialgebra-automorphism witness: $I_{\sigma_2} = \mathrm{Ad}_{U_W} \circ I_{\sigma_1}$.
\end{proof}

\subsubsection{What changes under wall-crossing, and what does not}

Five items:

\begin{enumerate}[leftmargin=2em]
 \item The BPS invariants $\Omega_\sigma(\gamma)$ \emph{change} chamber-by-chamber;
 this is the content of the wall-crossing formula.
 \item The tensor-factor ordering in the BPS Sym-decomposition of
 $\cH(\cC)$ \emph{changes};
 \item The isomorphism class of the bialgebra $\cH(\cC)$ \emph{does not change};
 \item The leading-residue bialgebra pairing $(-, -)_\cH|_{\mathrm{res}}$
 \emph{does not change} numerically (theorem);
 \item The all-mode bialgebra pairing is chamber-invariant under
 \cref{conj:hopf-pairing-all-mode-chamber-invariance}
 (conjecture).
\end{enumerate}

Items (3) and (4) state: the bialgebra $\cH(\cC)$ together with its
leading-residue pairing is a chamber-free invariant of the CY$_3$
category. Item (5) extends chamber-invariance of the pairing to all
modes conditional on
\cref{conj:hopf-pairing-all-mode-chamber-invariance}.
The KS wall-crossing cocycle is the explicit bialgebra-automorphism
witness that glues the chamber-by-chamber \emph{presentations} of
$\cH(\cC)$ into a single bialgebra; the Joyce--Song integration map
\cref{thm:js-bialgebra-chamber-equivariance} is the concrete chain-level
formula. The Drinfeld-double antipode is installed downstream
(\cref{rem:hopf-scope-stratified}), outside the chamber-gluing
datum.

\subsection{Stratum-independence}
\label{sec:wc-stratum-independence}

Unlike the gluing modes of Parts II--V (toric charts, McKay orbifolds,
derived-Morita tilting, factorisation-homology assembly), wall-crossing
gluing \emph{does not rely on the four-stratum taxonomy}. The KS quantum
torus and the Joyce--Song map are intrinsic to any CY$_3$ category with a
stability condition, regardless of whether the underlying geometry is:

\begin{itemize}[leftmargin=2em]
 \item toric (Stratum i): wall-crossing chambers stratify the K\"ahler
 cone, and the KS-$\Psi$-cocycle matches the GLY bond-factor cocycle
 of \cref{sec:gluing-p2-multi-chart};
 \item automorphism-equivariant (Stratum ii): the automorphism group acts
 on $\Stab(\cC)$ and permutes chambers; the chamber-quotient is
 preserved;
 \item orbifold-inertia (Stratum iii): each conjugacy-class component of
 the inertia stack has its own stability manifold, and wall-crossing
 within each component produces its own KS cocycle;
 \item lattice-polarised (Stratum iv): the Gritsenko--Borcherds-lift
 automorphic form produces generating functions of the same
 chamber-dependent BPS data; wall-crossing is a modular cocycle on
 the Kuga--Satake period domain.
\end{itemize}

In every case the \v{C}ech-cocycle structure is the same: chambers index
local charts; wall-crossing formulas provide transition maps; the BPS
spectrum is globally a section of the chamber sheaf.

\begin{remark}[Relationship to the motivic KS conjecture]
 The content of \cref{gl:thm:ks-wall-crossing} is one-half of the
 Kontsevich--Soibelman integrality conjecture; the other half, that the
 motivic BPS series factorises across the universal Pochhammer product,
 is now a theorem of Davison--Meinhardt in the critical-cohomology
 category. Together they give the full gluing picture: chamber-by-chamber
 local data, motivic wall-crossing cocycle, and global chamber-free
 algebra, all at the chain level.
\end{remark}

\subsection{Worked example: the conifold}
\label{sec:wc-conifold}

To demonstrate the formalism at chain level, we work out the wall-crossing
structure for the resolved conifold
$\bar{Y} = \Tot(\cO_{\mathbb P^1}(-1) \oplus \cO_{\mathbb P^1}(-1))$.

\subsubsection{Charge lattice and Euler form}

The derived category $D^b(\bar{Y})$ admits a Nagao--Nakajima / Szendr\H{o}i
NCCR Szendr\H{o}i (2008, arXiv:0705.3419) presenting it as
$D^b(\mathrm{mod}\text{-}\Jac(Q_{\mathrm{con}}, W_{\mathrm{con}}))$ where $Q_{\mathrm{con}}$ is the two-vertex
Klebanov--Witten quiver with arrows $a_{1,2}\colon 0 \to 1$,
$b_{1,2}\colon 1 \to 0$ and potential
$W_{\mathrm{con}} = a_1 b_1 a_2 b_2 - a_1 b_2 a_2 b_1$. The charge lattice is
$\Gamma = K_0(D^b(\bar{Y})) \cong \Z^2$ generated by simples
$\gamma_0 = [S_0]$, $\gamma_1 = [S_1]$ of the two vertices. The CY$_3$
Jacobi-algebra Ext-pattern between distinct simples reads
$\dim \Ext^0(S_i, S_j) = 0 = \dim \Ext^3(S_i, S_j)$ for $i \neq j$, while
$\dim \Ext^1(S_i, S_j) = \#\{\text{arrows } i \to j\} = 2$ and
$\dim \Ext^2(S_i, S_j) = \dim \Ext^1(S_j, S_i) = 2$ by CY$_3$ Serre
duality. The Euler pairing is therefore
\[
 \chi(\gamma_0, \gamma_1) \;=\; 0 - 2 + 2 - 0 \;=\; 0,
 \qquad \chi(\gamma_i, \gamma_i) \;=\; 1 - 0 + 0 - 1 \;=\; 0.
\]
The Euler form vanishes identically on $\Gamma$: the symmetric contribution
from $\Ext^1$ and $\Ext^2$ cancels, and the diagonal vanishes by CY$_3$
antisymmetry (\cref{prop:cy3-serre-antisymmetry}). The skew-Euler form is
$\langle \gamma_0, \gamma_1 \rangle = 0$ and the quantum torus generators
strictly commute,
\[
 \hat{x}_{\gamma_0} \hat{x}_{\gamma_1} \;=\; \hat{x}_{\gamma_0 + \gamma_1}
 \;=\; \hat{x}_{\gamma_1} \hat{x}_{\gamma_0},
\]
for the $K_0$-only charge lattice. Non-commutativity on the conifold enters
only when the lattice is enlarged to include the Chern-character refinements
$[\cO_{\mathrm{pt}}] = \gamma_0 + \gamma_1$ and $[\cO_{\mathbb P^1}(-1)]$
whose pairing with the simples generates the non-trivial commutators; the
motivic non-commutativity in the KS wall-crossing is carried by the
Chern-character refinement, not the simple-object lattice.

\subsubsection{Stability chambers and DT/PT generating functions}

The stability manifold $\Stab(D^b(\bar{Y}))$ is four-real-dimensional;
its chamber structure (Bridgeland Bridgeland (2011, arXiv:1002.4372),
Nagao--Nakajima Nagao--Nakajima (2011, arXiv:0809.2992)) is indexed by a discrete
family of NCCR tiltings related by mutations. Two canonical chambers:
\begin{itemize}[leftmargin=2em]
 \item \textbf{DT chamber}: stable bundle $\cO_{\bar{Y}} \to \bar{Y}$
 is a generator; generating function
 \[
 Z^{\mathrm{DT}}(\bar{Y})(q, Q) \;=\; M(-q)^{2} \prod_{k \geq 1}\big(1 - Q(-q)^k\big)^{k},
 \]
 where $M(q) = \prod_{k \geq 1}(1 - q^k)^{-k}$ is the MacMahon
 function (Maulik--Nekrasov--Okounkov--Pandharipande
 Maulik--Nekrasov--Okounkov--Pandharipande (2006, math/0312059, math/0406092)).
 \item \textbf{PT chamber}: stable Pandharipande--Thomas pairs;
 generating function
 $Z^{\mathrm{PT}}(\bar{Y})(q, Q) = \prod_{k \geq 1}(1 - Q(-q)^k)^{k}$, the
 DT/PT transform dropping the MacMahon prefactor $M(-q)^2$.
\end{itemize}
The DT/PT ratio $Z^{\mathrm{DT}}/Z^{\mathrm{PT}} = M(-q)^2$ is the Behrend-weighted
generating function of $0$-dimensional subschemes of $\bar{Y}$
(two copies of $\mathrm{Hilb}^\bullet(\C^3)$, one per conifold chart;
Bryan--Morrison (2009, arXiv:0812.5007) Thm.~1.1).

\subsubsection{Wall-crossing as pentagon identity}

The wall-crossing between DT and PT chambers is the motivic shadow of
the pentagon identity \cref{thm:psi-pentagon}. To make this explicit:
the DT chamber generates BPS states $\{n \gamma_0 + m \gamma_1 : n, m \geq 0\}$
with
$\Omega^{\mathrm{DT}}(\gamma_0) = \Omega^{\mathrm{DT}}(\gamma_1) = 1$,
$\Omega^{\mathrm{DT}}(\gamma_0 + \gamma_1) = -2$ (the two curve classes
$\cO_{\mathbb P^1}(-1)$), and $\Omega^{\mathrm{DT}}(k(\gamma_0 + \gamma_1)) = -2$
for all $k \geq 1$. The PT chamber has the mutated spectrum with the
same $\Omega(\gamma_0 + \gamma_1)$ but reorganised phase ordering.
The wall-crossing relation reads
\[
 \prod_{k \geq 1} \Psi(\hat{x}_{k(\gamma_0 + \gamma_1)})^{-2k}
 \cdot \Psi(\hat{x}_{\gamma_1})^{1} \cdot \Psi(\hat{x}_{\gamma_0})^{1}
 \;=\;
 \Psi(\hat{x}_{\gamma_0})^{1} \cdot \Psi(\hat{x}_{\gamma_1})^{1} \cdot \prod_{k \geq 1} \Psi(\hat{x}_{k(\gamma_0 + \gamma_1)})^{-2k},
\]
which restricts on the sublattice $\Z\gamma_0 \oplus \Z\gamma_1$ to a
repeated application of the pentagon \cref{thm:psi-pentagon}, one
per curve class $k(\gamma_0 + \gamma_1)$. Expanding both sides as
formal series in $\hat{x}_{\gamma_0}, \hat{x}_{\gamma_1}$ recovers
$Z^{\mathrm{DT}} = M(-q)^2 \cdot Z^{\mathrm{PT}}$, matching the MNOP
formula.

\subsubsection{Chamber-free CoHA and shifted quiver Yangian}

The chamber-free critical CoHA on the conifold is
\[
 \cH(\bar{Y}) \;\cong\; Y^+_{\hbar}(\widehat{\mathfrak{gl}}(1|1)),
\]
the shifted affine super-Yangian of $\mathfrak{gl}(1|1)$ in the GLY
normalisation Galakhov--Li--Yamazaki (2021, arXiv:2008.07006). The two supercharges
reflect the two simples $\gamma_0, \gamma_1$; the $\mathbb Z/2$ parity
records the two mutation-related tilting presentations. The chamber
dependence lives entirely in the shift parameter
$\zeta_\sigma \in \Gamma^\vee$, which parametrises which tilting
presentation realises the shifted Yangian at $\sigma$. The leading-residue
bialgebra pairing on
$Y^+_{\hbar}(\widehat{\mathfrak{gl}}(1|1))^{\mathrm{con}}$ takes
chamber-invariant numerical values by
\cref{thm:hopf-pairing-chamber-invariance} (theorem, bialgebra
scope; antipode reserved for the Drinfeld double
$D(Y^+) = Y_{\hbar}(\widehat{\mathfrak{gl}}(1|1))$); the all-mode propagation
of the asymmetric super-trace correction $c_{11} = -1$ is chamber-invariant
conditional on \cref{conj:hopf-pairing-all-mode-chamber-invariance}
. The DT and PT generating functions differ only by the
chamber-dependent prefactor $M(-q)^2$, which is the
``$\G_m$-equivariant zero-dimensional'' factor orthogonal to the
bialgebra pairing at leading residue.

\subsection{Summary and forward pointers}
\label{sec:wc-summary}

\begin{itemize}[leftmargin=2em]
 \item \cref{sec:ks-quantum-torus}: the quantum torus $\mathcal{T}_\Gamma$ encodes the
 Euler form on the charge lattice; on CY$_3$, Serre duality forces
 Euler antisymmetry, so the torus is motivically non-commutative with
 commutator $\LL^{2\chi(\gamma, \gamma')}$; in the classical limit it
 is the commutative group algebra of $\Gamma$.
 \item \cref{sec:chambers-gluing-cells}: chambers are local charts of
 $\Stab(\cC)$; walls are the codim-one boundaries; the
 quantum-dilogarithm cocycle $\prod \Psi(\hat{x}_\gamma)^{\Omega(\gamma)}$
 provides the \v{C}ech-gluing transition maps; the global BPS
 structure is a class in $\check{H}^1(\Stab(\cC), \Aut(\hat{\mathcal{T}}_\Gamma))$.
 \item \cref{sec:seiberg-mutation-gluing}: Seiberg duality / quiver
 mutation is the \emph{quiver-level} gluing automorphism; the
 mutation groupoid (a Weyl group for ADE) acts on the chamber
 lattice and lifts through Keller--Yang derived Morita
 (\cref{thm:keller-yang-mutation-wc}) to a group of bialgebra
 automorphisms $\mu_i^\cH \in \Aut_{\mathrm{bialg}}(\cH(\cC))$
 (\cref{thm:mutation-bialgebra-automorphism}); the CoHA is
 chamber-independent as a bialgebra, chamber data is carried by
 the shift parameter of the shifted quiver Yangian, and the
 conifold admits a three-chamber chain
 NCDT $\to$ DT $\to$ PT (\cref{thm:conifold-seiberg-chain})
 with conjectural $B_3$-braid lift
 (\cref{conj:braid-b3-conifold}).
 \item \cref{sec:js-hopf-chamber-independence}: the Joyce--Song
 integration map realises wall-crossing at the stack-function level,
 with the chamber-equivariance $I_{\sigma_2} = \mathrm{Ad}_{U_W} \circ
 I_{\sigma_1}$ and exponent the rational DT invariant
 $\bar{\Omega}_{\sigma_1}$. The critical CoHA is an associative
 bialgebra $(\mu, \iota, \Delta, \epsilon)$ and its leading-residue
 bialgebra pairing is preserved by every inner bialgebra automorphism
 $\mathrm{Ad}_{U_W}$ of $\cH(\cC)$ by $K_0$-orthogonality of the
 $\alpha_W$-corrections (theorem); all-mode chamber-invariance is
 conditional on
 \cref{conj:hopf-pairing-all-mode-chamber-invariance}. The
 Drinfeld-double antipode is installed downstream and is strict Hopf
 at rational/trigonometric/toroidal strata, quasi-Hopf with
 Felder--Jimbo--Konno dynamical associator at the elliptic stratum.
\end{itemize}

\bigskip
\noindent \textbf{Soibelman voice.} The content of this Part is the
motivic gluing cocycle: the BPS spectrum of a CY$_3$ category is not a
chamber-local datum but a \v{C}ech-cohomological class on $\Stab(\cC)$
taking values in $\Aut(\hat{\mathcal{T}}_\Gamma)$. Every chamber furnishes a
presentation of the same chamber-free algebra $\cH(\cC)$; the
wall-crossing automorphisms are the transition functions; the pentagon
identity is the local-model cocycle condition at a rank-two wall.
Integrality (Kontsevich--Soibelman 2008, Davison--Meinhardt 2020) is
the statement that this cohomology class descends to the $\G_m$-fixed
BPS Lie algebra $\mathfrak{g}_{\mathrm{BPS}}$, and its generating function is the
motivic refined DT invariant.

\S\ref{sec:gluing:lattice} takes up the lattice-polarised / automorphic gluing, which is the
gluing mode specifically adapted to K3-type geometries where the Borcherds
lift organises the chamber-free algebra into a modular form on the
Kuga--Satake period domain. The two gluing modes (wall-crossing and
automorphic) cooperate rather than compete on $K3 \times E$: wall-crossing
organises the BPS spectrum chamber-by-chamber on Stab, while the Borcherds
lift $\Delta_5$ organises the same spectrum globally as Fourier coefficients
of an automorphic form.

\bigskip
\paragraph{Primary sources cited in this Part.}
Kontsevich--Soibelman \cite{KS2008} \S 2.3--6.3 for the quantum torus, the
quantum dilogarithm, and the wall-crossing formula. Joyce--Song
Joyce--Song (2012, arXiv:0810.5645) Thm 5.8, \S 6.5 for the integration map $I$ and its
chamber-equivariance. Keller--Yang Keller--Yang (2011, arXiv:0906.0761) for
mutation-as-derived-equivalence. Galakhov--Li--Yamazaki
Galakhov--Li--Yamazaki (2021, arXiv:2008.07006) for shifted quiver Yangians indexing
chambers. Davison--Meinhardt Davison--Meinhardt (2020, arXiv:1601.02479) for the
chamber-independence of the critical CoHA. All statements in this Part are
stated at chain level; the upgrades to $(\infty, 1)$-categorical functoriality
under $\Phi$ are treated in the subsequent sections of this chapter.
